%% kapa: exam=1 round=23 year=2012 period=1 form=A subject=economics subject_ko=경제원론 seq=36 number=76
%% source: Q-Net 감정평가사 기출문제 (cst00309) artlSeq=5052578 file=raw/1008/5052578/2012년 제23회 감정평가사 제1차시험 기출문제(1교시 A형).pdf
%% conversion: kordoc 4.24.51 -> tools/gen1.py (auto); review=proofread-sonnet; flags=none
\begin{question}{76}
우리나라에서 산정되는 물가지수에 관한 설명으로 옳은 것은?
\begin{choices}
  \item 소비자물가지수 산정에 포함되는 재화와 용역은 해마다 달라진다.
  \item GDP 디플레이터 산정에는 파셰 지수(Paasche Index) 산식을 사용한다.
  \item 소비자물가지수 산정에는 국내에서 생산되는 재화와 용역만 포함된다.
  \item 생산자물가지수 산정에 포함되는 재화와 용역은 해마다 달라진다.
  \item GDP 디플레이터 산정에 포함되는 재화와 용역은 5년 마다 달라진다.
\end{choices}
\end{question}
\answer{2}
